\section{Documentación Detallada por Módulos}

\subsection{Módulo 1: Detección en Tiempo Real}

\subsubsection{Descripción General}

El módulo de detección proporciona visualización en tiempo real del procesamiento de video con YOLOv8, mostrando estadísticas actualizadas y permitiendo seleccionar entre múltiples cámaras.

\subsubsection{Componentes del Frontend}

\textbf{Archivo}: \texttt{resources/views/modulo1.blade.php} (323 líneas)

\paragraph{Estructura HTML Principal (Líneas 1-50)}

\begin{itemize}
    \item \textbf{Línea 1}: Extiende layout principal \texttt{@extends('layouts.app')}
    \item \textbf{Líneas 10-15}: Selector de cámara con ID \texttt{cameraSelector}
    \item \textbf{Línea 25}: Canvas para video \texttt{<canvas id="videoCanvas" width="800" height="600">}
    \item \textbf{Líneas 30-45}: Cards de estadísticas con IDs:
    \begin{itemize}
        \item \texttt{currentVehicles} - Vehículos actuales
        \item \texttt{totalVehicles} - Total detectados
        \item \texttt{fps} - Frames por segundo
        \item \texttt{averageVehicles} - Promedio de vehículos
    \end{itemize}
\end{itemize}

\paragraph{Lógica JavaScript (Líneas 150-323)}

\textbf{Variables Globales (Líneas 152-160)}:
\begin{lstlisting}[style=jsstyle, caption={Variables Globales - modulo1.blade.php}]
let websocket = null;
let currentCamera = 'CAM_002';
const canvas = document.getElementById('videoCanvas');
const ctx = canvas.getContext('2d');
let statsInterval = null;
let reconnectAttempts = 0;
const MAX_RECONNECT_ATTEMPTS = 5;
\end{lstlisting}

\textbf{Función connectWebSocket() (Líneas 170-210)}:
\begin{itemize}
    \item \textbf{Línea 171}: Construye URL WebSocket: \texttt{ws://localhost:8080/ws/stream?camera\_id=\$\{currentCamera\}}
    \item \textbf{Línea 173}: Crea instancia WebSocket con \texttt{new WebSocket(wsUrl)}
    \item \textbf{Línea 174}: Configura tipo binario: \texttt{websocket.binaryType = 'arraybuffer'}
    \item \textbf{Líneas 176-179}: Handler \texttt{onopen} - actualiza status a "Online"
    \item \textbf{Líneas 181-195}: Handler \texttt{onmessage}:
    \begin{itemize}
        \item Línea 182: Recibe ArrayBuffer del frame JPEG
        \item Línea 183: Convierte a Blob: \texttt{const blob = new Blob([event.data], \{type: 'image/jpeg'\})}
        \item Línea 184: Crea Object URL temporal
        \item Líneas 187-191: Dibuja imagen en canvas usando \texttt{ctx.drawImage()}
        \item Línea 192: Libera memoria con \texttt{URL.revokeObjectURL(url)}
    \end{itemize}
    \item \textbf{Líneas 197-204}: Manejo de errores y reconexión automática
\end{itemize}

\textbf{Función updateStats() (Líneas 215-250)}:
\begin{itemize}
    \item \textbf{Línea 217}: Fetch a API Laravel: \texttt{/api/vehicle-monitor/\$\{currentCamera\}}
    \item \textbf{Línea 218}: Parsea respuesta JSON
    \item \textbf{Líneas 220-223}: Actualiza elementos DOM:
    \begin{lstlisting}[style=jsstyle]
document.getElementById('currentVehicles').textContent = data.current_vehicles || 0;
document.getElementById('totalVehicles').textContent = data.total_detected || 0;
document.getElementById('fps').textContent = (data.fps || 0).toFixed(1);
    \end{lstlisting}
    \item \textbf{Líneas 225-230}: Llama \texttt{updateVehicleTypes(data.vehicle\_types)} para actualizar gráfico
    \item \textbf{Líneas 232-237}: Llama \texttt{updateVehicleColors(data.vehicle\_colors)} para badges de colores
\end{itemize}

\textbf{Inicialización (Líneas 300-315)}:
\begin{lstlisting}[style=jsstyle, caption={Inicialización del Módulo}]
window.addEventListener('DOMContentLoaded', () => {
    connectWebSocket();
    statsInterval = setInterval(updateStats, 1000); // Actualizar cada 1s
    
    document.getElementById('cameraSelector').addEventListener('change', (e) => {
        currentCamera = e.target.value;
        if (websocket) websocket.close();
        connectWebSocket();
    });
});
\end{lstlisting}

\subsubsection{Backend API}

\textbf{Archivo}: \texttt{app/Http/Controllers/VehicleMonitorController.php} (59 líneas)

\paragraph{Clase VehicleMonitorController}

\textbf{Propiedades (Líneas 12-13)}:
\begin{lstlisting}[style=phpstyle]
protected $detectorUrl = 'http://localhost:8080';
\end{lstlisting}

\textbf{Método getStats() (Líneas 20-57)}:
\begin{itemize}
    \item \textbf{Línea 21}: Recibe parámetro \texttt{\$cameraId} (default: 'CAM\_002')
    \item \textbf{Líneas 23-26}: Valida y normaliza cameraId:
    \begin{lstlisting}[style=phpstyle]
$cameraId = strtoupper($cameraId);
if (!in_array($cameraId, ['CAM_001', 'CAM_002'])) {
    $cameraId = 'CAM_002';
}
    \end{lstlisting}
    \item \textbf{Línea 28}: Construye URL completa: \texttt{"\{detectorUrl\}/api/vehicles?camera\_id=\{cameraId\}"}
    \item \textbf{Línea 29}: Hace petición HTTP con timeout de 5s: \texttt{Http::timeout(5)->get(\$url)}
    \item \textbf{Líneas 31-33}: Si exitoso, retorna JSON del detector Python
    \item \textbf{Líneas 35-42}: Si falla, retorna estructura por defecto con valores en 0
\end{itemize}

\subsubsection{Servicio Python Detector}

\textbf{Archivo}: \texttt{deteccion/vehiculo\_detector.py} (767 líneas)

\paragraph{API Flask - Endpoint /api/vehicles (Líneas 550-580)}

\textbf{Función get\_vehicle\_stats() (Línea 552)}:
\begin{lstlisting}[style=pythonstyle, caption={Endpoint de Estadísticas}]
@app.route('/api/vehicles', methods=['GET'])
def get_vehicle_stats():
    camera_id = request.args.get('camera_id', 'CAM_002')
    
    with data_lock:
        state = camera_states.get(camera_id, {})
        tracked = state.get('tracked_vehicles', {})
        
        # Contar por tipo de vehículo
        vehicle_types = {}
        vehicle_colors = {}
        
        for vehicle_id, vehicle_data in tracked.items():
            v_type = vehicle_data.get('type', 'desconocido')
            v_color = vehicle_data.get('color', 'desconocido')
            
            vehicle_types[v_type] = vehicle_types.get(v_type, 0) + 1
            vehicle_colors[v_color] = vehicle_colors.get(v_color, 0) + 1
        
        return jsonify({
            'camera_id': camera_id,
            'status': state.get('status', 'offline'),
            'current_vehicles': len(tracked),
            'total_detected': state.get('total_vehicles_detected', 0),
            'fps': state.get('fps', 0.0),
            'vehicle_types': vehicle_types,
            'vehicle_colors': vehicle_colors
        })
\end{lstlisting}

\paragraph{WebSocket Streaming - Endpoint /ws/stream (Líneas 600-650)}

\textbf{Función stream\_video() (Línea 602)}:
\begin{itemize}
    \item \textbf{Línea 604}: Obtiene \texttt{camera\_id} de query params
    \item \textbf{Líneas 608-615}: Loop infinito que:
    \begin{itemize}
        \item Línea 610: Accede a \texttt{camera\_states[camera\_id]['processed\_frame']} con lock
        \item Línea 612: Si frame existe, lo codifica a JPEG usando \texttt{cv2.imencode('.jpg', frame)}
        \item Línea 614: Envía bytes vía WebSocket: \texttt{ws.send(frame\_bytes)}
        \item Línea 616: Espera 33ms (30 FPS): \texttt{time.sleep(0.033)}
    \end{itemize}
    \item \textbf{Líneas 620-625}: Manejo de excepciones y cierre de conexión
\end{itemize}

\paragraph{Worker Thread - worker\_skyline() (Líneas 350-450)}

\textbf{Función principal (Línea 352)}:
\begin{itemize}
    \item \textbf{Líneas 355-365}: Llama \texttt{get\_skyline\_stream\_url\_robust()} para obtener URL HLS
    \item \textbf{Línea 370}: Abre stream con OpenCV: \texttt{cap = cv2.VideoCapture(stream\_url)}
    \item \textbf{Líneas 375-420}: Loop de procesamiento:
    \begin{itemize}
        \item Línea 377: Lee frame: \texttt{ret, frame = cap.read()}
        \item Línea 380: Llama \texttt{process\_frame\_generic(frame, 'CAM\_002')}
        \item Líneas 385-395: Calcula FPS usando \texttt{time.time()} y contador de frames
    \end{itemize}
\end{itemize}

\textbf{Función process\_frame\_generic() (Líneas 200-280)}:
\begin{itemize}
    \item \textbf{Línea 205}: Ejecuta YOLO: \texttt{results = model(frame, conf=0.4)}
    \item \textbf{Líneas 210-250}: Para cada detección:
    \begin{itemize}
        \item Línea 215: Obtiene bbox, confianza y clase
        \item Línea 220: Filtra solo vehículos (car, motorcycle, bus, truck)
        \item Línea 225: Llama \texttt{match\_vehicle\_to\_tracked(bbox, tracked, max\_distance=50)}
        \item Líneas 230-240: Si es nuevo vehículo:
        \begin{itemize}
            \item Asigna nuevo ID: \texttt{vehicle\_id = state['next\_vehicle\_id']}
            \item Incrementa contador
            \item Llama \texttt{get\_dominant\_color(frame, bbox)}
            \item Llama \texttt{save\_detection\_to\_xml(vehicle\_data, camera\_id)}
        \end{itemize}
        \item Línea 245: Dibuja bbox en frame con \texttt{cv2.rectangle()}
    \end{itemize}
    \item \textbf{Línea 270}: Almacena frame procesado: \texttt{state['processed\_frame'] = frame}
\end{itemize}

\subsubsection{Flujo Completo de Detección}

\begin{figure}[H]
\centering
\begin{tikzpicture}[
    box/.style={rectangle, draw, fill=cyan!15, minimum width=3.5cm, minimum height=0.7cm, font=\small},
    node distance=0.8cm
]
    \node[box] (step1) {1. Usuario abre modulo1.blade.php};
    \node[box, below=of step1] (step2) {2. DOMContentLoaded ejecuta (L.300)};
    \node[box, below=of step2] (step3) {3. connectWebSocket() (L.170)};
    \node[box, below=of step3] (step4) {4. WebSocket a localhost:8080/ws/stream};
    \node[box, below=of step4] (step5) {5. Python stream\_video() envía JPEGs};
    \node[box, below=of step5] (step6) {6. onmessage dibuja en canvas (L.181)};
    \node[box, below=of step6] (step7) {7. setInterval llama updateStats() (L.305)};
    \node[box, below=of step7] (step8) {8. Fetch a /api/vehicle-monitor};
    \node[box, below=of step8] (step9) {9. VehicleMonitorController::getStats()};
    \node[box, below=of step9] (step10) {10. HTTP a Python /api/vehicles};
    \node[box, below=of step10] (step11) {11. Retorna vehicle\_types, colors, FPS};
    \node[box, below=of step11] (step12) {12. Actualiza DOM (L.220-237)};
    
    \draw[-{Latex}] (step1) -- (step2);
    \draw[-{Latex}] (step2) -- (step3);
    \draw[-{Latex}] (step3) -- (step4);
    \draw[-{Latex}] (step4) -- (step5);
    \draw[-{Latex}] (step5) -- (step6);
    \draw[-{Latex}] (step6) -- (step7);
    \draw[-{Latex}] (step7) -- (step8);
    \draw[-{Latex}] (step8) -- (step9);
    \draw[-{Latex}] (step9) -- (step10);
    \draw[-{Latex}] (step10) -- (step11);
    \draw[-{Latex}] (step11) -- (step12);
\end{tikzpicture}
\caption{Flujo Completo del Módulo Detección}
\end{figure}

\newpage
\subsection{Módulo 2: Reportes y Estadísticas}

\subsubsection{Descripción General}

El módulo de reportes permite visualizar el historial de detecciones con filtros avanzados, estadísticas agregadas y exportación a Excel/CSV.

\subsubsection{Controlador Principal}

\textbf{Archivo}: `texttt{app/Http/Controllers/ReporteController.php} (171 líneas)

\paragraph{Método obtenerRegistrosFiltrados() (Líneas 12-73)}

\textbf{Variables clave}:
\begin{itemize}
    \item \textbf{Línea 14}: `texttt{$registros = []} - Array acumulador de registros filtrados
    \item \textbf{Línea 15}: `texttt{$xmlPath = storage\_path('app/vehiculos\_db.xml')} - Ruta del XML
    \item \textbf{Línea 18-19}: Carga XML con `texttt{simplexml\_load\_string()}
\end{itemize}

\textbf{Bucle de filtrado (Líneas 22-65)}:
\begin{lstlisting}[style=phpstyle]
foreach (->deteccion as ) {
     = [
        'fecha' => (string)->fecha,
        'tipo'  => (string)->tipo,
        'confianza' => (float)->confianza,
        'color' => isset(->color) ? (string)->color : 'desconocido',
        'camara' => isset(->camara) ? (string)->camara : 'CAM_002'
    ];
    
     = true;
    
    // L.34-38: Filtro por cámara
    if (->has('camera_id') && ->camera_id != '') {
        if (['camara'] != ->camera_id) {
             = false;
        }
    }
    
    // L.40-44: Filtro por tipo (Auto, Moto, Bus, Camión)
    if (->has('tipo') && ->tipo != '') {...}
    
    // L.46-52: Filtro por fecha inicio
    if (->has('fecha_inicio') && ->fecha_inicio != '') {
         = strtotime(['fecha']);
         = strtotime(->fecha_inicio . ' 00:00:00');
        if ( < )  = false;
    }
    
    // L.54-60: Filtro por fecha fin
    if (->has('fecha_fin') && ->fecha_fin != '') {...}
    
    if () [] = ;
}
\end{lstlisting}

\textbf{Ordenamiento (Líneas 68-70)}: Ordena desc por fecha usando `texttt{usort()} y `texttt{strtotime()}

\paragraph{Método exportarExcel() (Líneas 82-140)}

\textbf{Variables clave (Líneas 84-88)}:
\begin{itemize}
    \item `texttt{$registros} - Obtiene registros filtrados
    \item `texttt{$nombreArchivo} - Genera nombre con timestamp: `texttt{reportes\_ . date('Y-m-d\_His') . '.csv'}
    \item `texttt{$rutaTemporal} - Path temporal: `texttt{storage\_path('app/temp')}
\end{itemize}

\textbf{Generación CSV (Líneas 95-125)}:
\begin{lstlisting}[style=phpstyle]
// L.100: Abrir archivo en modo escritura
 = fopen(, 'w');

// L.102-104: BOM UTF-8 para compatibilidad Excel
fprintf(, chr(0xEF).chr(0xBB).chr(0xBF));

// L.106-108: Escribir encabezados
fputcsv(, [
    'Fecha/Hora', 'Tipo Vehículo', 'Confianza (%)', 
    'Color', 'Cámara', 'Nombre Cámara'
], ',');

// L.110-120: Escribir datos
foreach ( as ) {
    fputcsv(, [
        ['fecha'],
        ['tipo'],
        number_format(['confianza'], 2),
        ['color'],
        ['camara'],
        ['nombre_camara']
    ], ',');
}

fclose(); // L.122
\end{lstlisting}

\textbf{Descarga y limpieza (Líneas 125-138)}:
\begin{itemize}
    \item \textbf{Línea 127}: `texttt{return response()->download($rutaCompleta, $nombreArchivo)->deleteFileAfterSend(true)}
    \item El método `texttt{deleteFileAfterSend(true)} elimina automáticamente el archivo temporal
\end{itemize}

